前两讲分别从力和能量两个角度考察了电场，引入了描述电场的矢量场
$\boldsymbol E$ 与标量场 $\varphi$。电场强度可以用于一般电磁问题，电势则专用于静电场。电场线和
等势面虽然是图形化的表达，却能把许多空间关系直接呈现出来。此前主要讨论电场对少数点电荷的作用，
本讲开始研究电场与宏观物质的相互作用。

所有宏观物质都由原子构成，原子由带正电的原子核与核外电子组成。把物质置于外电场中，正、负电荷会
受到方向相反的作用；物质内部电荷能否移动，决定了它对外电场的响应方式。

\subsection*{导体和绝缘体}

孤立原子中的电子受原子核束缚。许多原子聚集成固体以后，每个电子不仅受到最近原子核的作用，也受到
周围大量原子的影响。某些距离原子核较远的电子不再局限于单个原子，而可以在整个固体中运动。准确描述
这一现象需要量子理论；在宏观静电学中，可以把导体理解为含有大量可移动载流子的物质，把绝缘体理解为
电荷在外场作用下仍基本束缚在原处的物质。

金属是典型导体。向金属传递一些电荷后，电荷不会固定在接触点附近，而会在极短时间内重新分布到导体
表面。要让导体长期带电，必须用绝缘体把它与其他导电区域隔开，否则电荷会迅速流失。玻璃、橡胶和干燥
木材是常见绝缘体；它们经过摩擦后，电荷可以长时间停留在局部区域。

上述“自由电子”和“完全束缚电子”的图像是简化模型。真实材料的导电性还随温度、杂质、场强和频率
变化；绝缘体在足够强的电场下也会击穿。但这一模型已经足以建立静电平衡的主要规律。

应当进一步说明，这里所谓“自由”并不是电子完全不受原子核和晶格作用。电子在固体中始终处于所有离子
实产生的周期势场里，还会与晶格振动、杂质和其他电子碰撞。导体与绝缘体的本质差别，是材料中是否存在
能够在很小外场下改变占据状态并形成宏观电流的载流子。金属中的载流子数目很大，所以外场一出现，表面
电荷就能在极短时间内完成重新分布；绝缘体中的电荷只能产生很小的束缚位移，通常表现为极化而不是持续
导电。半导体处在二者之间，其载流子数目可被温度、光照或掺杂显著改变。

导体得到的多余电荷会遍布它的可达表面，而不会停在最初接触的位置。这里的“表面”包括外表面，也可能
包括含有内部电荷的空腔表面；究竟分布到哪一部分，要由静电平衡条件共同决定。若导体与另一个导体直接
接触，它们将构成一个新的整体等势体，电荷继续移动，直到整个连通导体达到共同电势。因此实验中常用
绝缘支架、绝缘手柄或真空把带电导体与环境隔离。

\subsection*{静电平衡}

把导体放入预先存在的外电场，内部载流子立即受力并发生定向移动。载流子到达导体边界后不能继续离开，
于是表面出现净电荷。这些表面电荷又产生电场，并与外电场叠加。只要导体内部总电场不为零，电荷分布就
还会继续改变；最终导体内部处处无场，宏观状态不再改变，这就是\textbf{静电平衡}。移动后留在表面的净
电荷称为\textbf{感应电荷}。

图3.1以一个带正电导体和一个净电荷为零的导体为例。中性导体靠近带电体的一侧出现负感应电荷，远侧
出现正感应电荷；电场线从正电荷出发，终止于负电荷或延伸到无穷远。导体表面本身是等势面，导体内部
没有电场线。

\conductorFieldFigure
\diagramcaption{图 3.1：静电平衡状态下两个导体的电场线、等势线与感应电荷分布。以两球为例，实线为场线、虚线为等势线；电场线垂直球面，延伸到画框外的线并非在框边终止。}

根据形成静电平衡的过程，可以得到以下性质。

\begin{enumerate}[label=\arabic*.]
  \item \textbf{导体材料内部电场处处为零。}如果某处 $\boldsymbol E\ne0$，可移动电荷仍会受力，电荷分布
  还会变化，因而不可能处于静电平衡。

  \item \textbf{导体材料内部没有净体电荷。}在导体内部取任意闭合 Gauss 面，因为面上
  $\boldsymbol E=0$，电通量为零，所包围的净电荷也必为零。因此导体的多余电荷只能分布在表面。

  \item \textbf{整个导体是等势体，表面是等势面。}导体内部任意两点间的电势差为
  $-\int\boldsymbol E\cdot\dd\boldsymbol l=0$。把导体连接到大地称为接地；大地可吸收或提供大量电荷，
  通常规定接地导体的电势为零。

  \item \textbf{导体表面外侧的电场垂直于表面。}若存在切向分量，表面电荷会沿导体表面流动。取一个
  横跨导体表面的极薄柱形 Gauss 面，导体内侧电场为零，得到表面外侧法向电场与面电荷密度的关系
  \begin{equation}
    \boxed{E_n=\frac{\sigma}{\varepsilon_0}}.
  \end{equation}
  其中 $E_n$ 以导体外法线方向为正。$\sigma>0$ 时电场指向导体外，$\sigma<0$ 时指向导体内。
\end{enumerate}

图3.2给出用静电感应使导体带正电的完整顺序。起初导体呈电中性；把带负电物体靠近导体，导体左侧聚集
正电荷、右侧聚集负电荷；保持负电体不动并将导体接地，电子受排斥流入大地，导体上只剩净正电荷；先
断开地线，导体上的正电荷被保留下来；最后移走负电体，正电荷重新分布到整个外表面。顺序不能颠倒：若
先移走带电体再断开地线，电荷会经地线回到大地，导体仍为中性。

\inductionSequenceFigure
\diagramcaption{图 3.2：一个完整的静电感应带电过程。图中负电棒始终位于导体左侧，接地只发生在步骤 (c)。}

感应电荷的总量和局部分布必须分开理解。未接地的中性导体在感应过程中总电荷始终为零，但近端、远端可
分别出现等量异号电荷；接地以后，地线允许电荷进入或离开，导体总电荷不再固定。断开地线只固定导体的
总电荷，并不固定表面电荷分布；当外部带电体继续移动时，表面电荷仍会沿导体重新排列。正因如此，静电
感应问题不能只画几个“正、负号”就结束，还必须同时检查总电荷守恒、导体电势条件和表面法向电场方向。

接地也不能简单理解为“把所有电荷都导走”。它真正固定的是导体电势。为了使电势达到零，大地可能吸收
电子，也可能向导体提供电子，方向取决于外部电荷的符号和几何位置。图3.2中负电棒排斥电子，所以电子
从导体流向大地；若靠近的是正电棒，过程中的电荷流向和最后带电符号都将反转。

静电平衡时，同一导体表面不可能有一条电场线从正电荷区域出发、再终止于该导体的负电荷区域。因为沿
电场线电势必须下降，而同一导体各点电势相等。把这一结论用于内部无电荷的封闭导体壳：空腔内表面任意
两点之间不能由电场线连接；又由包围空腔的 Gauss 面可知内表面总电荷为零，于是空腔内表面处处无电荷，
空腔内电场为零。这就是静电屏蔽。

法拉第笼是静电屏蔽的典型应用。完整金属外壳可以阻止外部静电场进入内部；网状金属笼对变化较慢、波长
远大于网孔的电磁场也有明显屏蔽作用。屏蔽结构可保护敏感仪器免受射频干扰，也可以把无线电发射设备的
干扰限制在内部。遭遇雷击或静电放电时，电流主要沿导体笼外表面传导，内部人员和设备受到保护。

图 3.3 的金属网构成连通屏蔽壳；在上述静电条件下，壳内无电荷的空腔与外界电场隔离。

\faradayCageFigure
\diagramcaption{图 3.3：法拉第笼的屏蔽情景。理想静电极限下腔内无场；雷击瞬态的安全还依赖连接、接地及屏蔽完整性。}


% EM teaching enrichment 20260927
\paragraph{带电空腔：把内表面与外表面分开}
中性孤立导体壳内放入电荷 $q$，在导体材料内取包围空腔的高斯面，因 $E=0$，有
$q+Q_{\rm inner}=0$。再由壳的净电荷守恒得 $Q_{\rm outer}=q$。这两步只决定总量；内电荷偏心时内表面电荷一般不均匀。
若外表面为球面、外界无其他电荷，外部解由球形等势边界确定为球心点电荷 $q$ 的场，与空腔内电荷偏心方向无关。
若壳接地且外界无其他电荷，外表面可以不带电，外部场为零；此时壳总电荷已由地线改变，不能再使用原来的中性约束。

\subsection*{唯一性定理}

定性规律可以判断静电平衡的许多特征；要确定空间中电势和电场的具体分布，还需要静电场的唯一性定理。
由
$\boldsymbol E=-\boldsymbol\nabla\varphi$ 与 $\boldsymbol\nabla\cdot\boldsymbol E=\rho/\varepsilon_0$
联立，得到
\begin{equation}
  \boxed{\frac{\partial^2\varphi}{\partial x^2}
  +\frac{\partial^2\varphi}{\partial y^2}
  +\frac{\partial^2\varphi}{\partial z^2}
  =-\frac{\rho}{\varepsilon_0}}.
\end{equation}
这就是 Poisson 方程。无电荷区域中 $\rho=0$，它退化为 Laplace 方程
$\nabla^2\varphi=0$。

中学阶段不需要直接求解一般的三维偏微分方程，但离散化可以直观说明为什么必须给出边界条件。设电荷和
电势只随 $x$ 变化，则
\begin{equation}
  \boxed{\frac{\dd^2\varphi}{\dd x^2}=-\frac{\rho(x)}{\varepsilon_0}}.
\end{equation}

离散化的思想与测量城市温度很相似。温度本来是位置的连续函数，但实际不可能在每一个数学点都放置温度
计；只要在足够密的观测站取得温度，再利用相邻点之间的变化，就能近似恢复整座城市的温度分布。电势也
一样：先在有限个节点上求 $\varphi_n$，节点越密，离散结果越接近连续函数。Poisson 方程的离散式不是
额外的物理假设，而是连续二阶导数的差分近似。
在 $A,B$ 之间取间距均为 $a$ 的离散点，点 $n$ 处的电势为 $\varphi_n$。其右、左两侧一阶导数分别近似为
$\bigl(\varphi_{n+1}-\varphi_n\bigr)/a$ 与
$\bigl(\varphi_n-\varphi_{n-1}\bigr)/a$，因此二阶导数的离散形式为
\begin{equation}
  \boxed{\frac{\varphi_{n+1}+\varphi_{n-1}-2\varphi_n}{a^2}
  =-\frac{\rho_n}{\varepsilon_0}}.
\end{equation}

图 3.4 中节点 $n$ 与两侧相邻节点均相距 $a$，所以其二阶差分恰由上式三项构成；$A,B$ 的电势不再作为未知量。

\discretePotentialFigure
\diagramcaption{图 3.4：把 $A,B$ 之间连续的电势分布离散化。相邻节点间距均为 $a$，端点电势是边界条件。}

对每一个内部节点都列一个方程，就得到包含大量未知电势的线性方程组。内部节点的电荷密度必须已知；对
最靠近两端的节点，方程还会包含端点 $A,B$ 的电势，因此边界电势也必须给定。令 $a\to0$，离散点无限
细密，便回到连续的 Poisson 方程。

静电场唯一性定理可以表述为：给定区域内的电荷分布，并在边界上给定适当条件，区域内的电势和电场便
唯一确定。常见边界条件有三类：边界电势已知；边界法向电场已知；一部分边界给电势、另一部分边界给
法向电场。导体是等势体，所以“某个导体电势已知”本身就是一种边界条件；若只给导体的总电荷，还需用
总电荷约束共同确定导体电势。

为什么内部电荷和边界条件缺一不可，可以直接从差分方程看出。对 $n=2,3,\ldots,N-2$，方程右侧由
$\rho_n$ 给定；但对最靠近左端的 $n=1$，方程包含端点未知量 $\varphi_0$，对最靠近右端的
$n=N-1$，方程包含 $\varphi_N$。若不给 $A,B$ 的电势，未知数就比独立方程多，整体可以发生不受限制的
改变。三维问题虽然节点更多、邻接关系更复杂，逻辑完全相同：内部的 Poisson 方程规定局部曲率，边界
条件负责把所有局部关系闭合成唯一的全局解。

若给定的是边界法向场而非边界电势，由 Gauss 定律还必须满足总通量与区域内总电荷相容。纯 Neumann
边界条件只决定电势到一个任意加法常量，但这个常量不影响电场；混合边界条件则把两类信息分别施加在
不同边界部分。这三种分支条件，正是 Dirichlet、Neumann 和混合边界条件在静电学中的对应。

唯一性定理的力量在于：只要构造出一个满足相同区域电荷分布和相同边界条件的电势函数，就不必再去求
复杂的感应电荷分布；这个构造出来的解就是唯一的物理解。

\subsection*{电像法}

静电平衡问题通常同时包含给定电荷和导体表面上未知的感应电荷。直接积分感应电荷非常困难。若能在所
研究区域之外放置若干假想点电荷，使它们与真实电荷共同产生的电势恰好满足导体边界条件，那么由唯一性
定理可知，假想电荷在研究区域内的作用与真实感应电荷完全等效。这就是\textbf{电像法}。像电荷不是真实
电荷，只能放在研究区域之外，不能把它当成物体的一部分去计算受力。

电像法的经典例子是：半径为 $R$ 的接地导体球，球心为 $C$；球外轴线上距球心 $L>R$ 处放一点电荷
$+Q$。在球内同一轴线上、距球心 $l$ 处放一个像电荷 $-q$。只要整个球面电势为零，球外电场就与真实
问题完全相同。取轴线上球面的近、远两个点 $a,b$，有
\[
  \frac{1}{4\pi\varepsilon_0}
  \left(\frac{Q}{L-R}-\frac{q}{R-l}\right)=0,
  \qquad
  \frac{1}{4\pi\varepsilon_0}
  \left(\frac{Q}{L+R}-\frac{q}{R+l}\right)=0.
\]
解得
\begin{equation}
  \boxed{q=\frac RLQ,\qquad l=\frac RL R=\frac{R^2}{L}}.
\end{equation}

这里只在球面上选取两个点列方程，看起来似乎不足以保证整个球面电势都为零。关键在于所得参数同时满足
$Ll=R^2$ 和 $q/Q=R/L$。对球面任意点 $P$，三角形几何给出
$r_1/r_2=L/R$，从而 $Q/r_1=q/r_2$；所以任意 $P$ 的电势都为零。近、远两点的方程只是求出候选参数，
任意点的距离比关系才完成对整个边界条件的验证。电像法中的每一次构造都应做这样的全边界检查，不能只
验证图上少数特殊点。

William Thomson（后来获封 Kelvin 勋爵）系统发展了这一方法。它之所以巧妙，是因为我们不必知道每一
小块球面上到底聚集了多少感应电荷，只需找到一个位于禁入区域内的简单点电荷系统，使球外电势满足完全
相同的方程和边界条件。像电荷的位置落在导体内部，那里不是所求球外区域，因此像电荷产生的奇点不会
污染真实解。

图 3.5 中 $C$ 是球心，$P$ 严格位于球面；$CP=R$，真实电荷与像电荷到 $C$ 的距离分别为 $L,l$。

\groundedSphereImageFigure
\diagramcaption{图 3.5：接地导体球的电像法几何。真实电荷为 $+Q$，像电荷为 $-q$，且 $Ll=R^2$。}

对球面上任意点 $P$，几何关系保证
$r_1/r_2=L/R=Q/q$，所以 $Q/r_1-q/r_2=0$；球面确为零等势面。像电荷的
位置和大小确定后，外电荷受到的吸引力等于它与像电荷之间的 Coulomb 力：
\[
  F=\frac{1}{4\pi\varepsilon_0}\frac{Qq}{(L-l)^2}.
\]
外电荷与接地球组成系统的静电势能不是两个真实点电荷势能，而是
\[
  U=\frac12Q\varphi_{\rm ind}(Q)
  =-\frac12\frac{1}{4\pi\varepsilon_0}\frac{Qq}{L-l}.
\]
因子 $1/2$ 来自导体感应电荷随真实电荷逐步建立；若把像电荷误当真实电荷，这个因子会被遗漏。

把真实电荷与像电荷的场在球外叠加，在球面处取法向分量，再用 $\sigma=\varepsilon_0 E_n$，得到距真实电荷为 $r_1$ 的球面
点上的感应电荷密度
\begin{equation}
  \boxed{\sigma=-\frac{1}{4\pi}\frac{Q(L^2-R^2)}{Rr_1^3}}.
\end{equation}
这里已经约去 $\varepsilon_0$；结果的单位为电荷面密度。球面各处 $\sigma<0$，且离 $Q$ 越近，$r_1$
越小，负感应电荷密度的绝对值越大。


% EM teaching enrichment 20260927
\paragraph{先用平面模型检验力与能量}
电荷 $q$ 距无限接地平面为 $a$，像电荷为镜像位置的 $-q$，二者相距 $2a$。
真实电荷受力大小 $kq^2/(4a^2)$，指向平面；系统与感应有关的势能却为
\[
 U(a)=\frac12q\varphi_{\rm ind}=-\frac{kq^2}{4a},\qquad
 F_a=-\frac{\dd U}{\dd a}=-\frac{kq^2}{4a^2}.
\]
求力时像电荷场可直接使用；求位置导数时像的位置和感应分布随 $a$ 一起改变。像法只在导体外替代感应场，像电荷所在处不是一个真实可访问的受力粒子。

\subsection*{等离子体频率}

导体最终达到静电平衡，不仅因为内部存在可移动载流子，还因为载流子的运动会耗散能量。竖直弹簧下突然
挂上重物，若没有阻力，重物不会停在新的静力平衡位置，而会围绕它持续振动；只有阻尼耗散机械能，重物
才会最终静止。导体中的载流子也一样：若完全没有碰撞和电阻，它们不会停在屏蔽外场的电荷分布上，而会
围绕电中性位置集体振荡。

把弹簧的过程说得更具体些：重物刚挂上时，原来的弹簧长度不是新的静力平衡位置。重力先使重物向下加速，
弹簧伸长后恢复力增大；重物经过平衡位置时仍有速度，于是继续向下，随后被拉回。没有空气阻力和弹簧内耗
时，这一往复过程永不停止。导体中载流子的惯性对应重物质量，表面电荷产生的恢复场对应弹簧恢复力，碰撞
和电阻则对应阻尼。通常金属中的碰撞时间很短，因此宏观观察只看到迅速建立的静电平衡，而看不到长时间
的自由振荡。

设导体中可移动载流子的质量为 $m$、电荷为 $q$、数密度为 $n$。外电场 $E_0$ 沿 $+x$ 方向突然加上，
载流子整体位移 $x$ 后，样品两侧出现面电荷，大小为 $\sigma=nqx$。这些面电荷在内部产生恢复场
\[
  E'=-\frac{nq}{\varepsilon_0}x.
\]

图中以正载流子向右位移表示正负面电荷的几何关系。
若载流子是电子，$q=-e$，在向右外场下平衡位移 $x_0<0$，电子向左；
两端面的符号仍应从乘积 $nqx$ 判断，不能仅由位移方向猜测。
\plasmaSlabFigure

忽略碰撞和阻尼，每个载流子的运动方程为
\[
  qE_0-\frac{nq^2}{\varepsilon_0}x=m\ddot x.
\]
整理为
\begin{equation}
  \boxed{\ddot x+\frac{nq^2}{m\varepsilon_0}x=\frac{qE_0}{m}}.
\end{equation}
这是围绕平衡位置 $x_0=\varepsilon_0E_0/(nq)$ 的简谐振动，固有角频率为
\begin{equation}
  \boxed{\omega_p=\left(\frac{nq^2}{m\varepsilon_0}\right)^{1/2}}.
\end{equation}

若平均阻力为 $-m\nu\dot x$，令 $y=x-x_0$，则
\[
 \ddot y+\nu\dot y+\omega_p^2y=0,\qquad
 s_\pm=\frac{-\nu\pm\sqrt{\nu^2-4\omega_p^2}}2.
\]
当 $0<\nu<2\omega_p$ 时为衰减振荡，振幅包络为 $e^{-\nu t/2}$；
$\nu=2\omega_p$ 时为临界阻尼；$\nu>2\omega_p$ 时为非振荡的过阻尼趋近。
特别是 $\nu\gg\omega_p$ 时，慢根 $s_+\simeq-\omega_p^2/\nu$，
弛豫时间约为 $\nu/\omega_p^2$，所以不能笼统说“碰撞越强，平衡越快”。
弱阻尼而外场交变时，在适当驱动下可于 $\omega\approx\omega_p$ 附近出现强响应。因此“静电平衡”和“等离子体振荡”并不矛盾：前者是有耗散、长时间极限，后者是惯性
不可忽略时的动力学行为。

等离子体是离子和电子组成、能够表现出集体行为的带电粒子体系。离子质量远大于电子，最简单近似下可把
离子视为静止背景，只考虑电子相对离子的振荡；此时带符号电荷为 $q=-e$、$m=m_e$，频率中的 $q^2=e^2$，$\omega_p$ 主要由电子数密度决定。
电子偏离电中性位置后，电荷分离产生恢复场；电子越过平衡位置时仍有惯性，于是形成等离子体振荡。

等离子体频率公式还说明等离子体频率只由载流子数密度、电荷量和质量决定，与样品在理想无限均匀近似下的横向
尺寸无关。电子的质量很小，所以电子等离子体频率通常很高；离子也可以振荡，但其特征频率低得多。在
有限等离子体中还会出现边界效应、温度效应和不同传播方向的多种波模，等离子体频率公式给出的是最基本的均匀纵向
集体振荡尺度。

地球约 $60\,\mathrm{km}$ 以上的高层大气受到太阳高能辐射和宇宙线作用而部分电离，形成电离层。入射电磁波的频率低于电离层的等离
子体频率时，电子能够集体响应并使电磁波被反射；频率足够高时，电磁波可以穿过电离层。短波通信可借助
电离层一次或多次反射跨越地平线，而卫星通信需要采用能够穿透电离层的较高频率。

图 3.6 中橙色路径表示被电离层折回的通信波，蓝色路径穿过电离层到达卫星；两条路径对应不同的频率条件。

\ionosphereFigure
\diagramcaption{图 3.6：电离层反射实现远距离通信。图中同时区分低频反射与高频穿透。}

现实导体中的碰撞使载流子振荡衰减，系统才会趋于静电平衡；在高频电磁场或等离子体中，惯性和集体振荡
则不能忽略。静电平衡与等离子体频率看似分属静力学和动力学，实际是同一套电荷重新分布机制在有、无
耗散两种极限下的表现。

% EM teaching enrichment 20260927
\paragraph{把振荡模型接回电荷弛豫}
在强碰撞极限，前面的电导模型给出 $\sigma=nq^2/(m\nu)$，从而
\[
 \frac{\nu}{\omega_p^2}=\frac{\varepsilon_0}{\sigma}.
\]
这正是忽略惯性后，由电荷连续性与 $\boldsymbol j=\sigma\boldsymbol E$ 得到的体电荷弛豫时间：
$\partial_t\rho=-\nabla\cdot\boldsymbol j=-\sigma\rho/\varepsilon_0$。
二者相符须满足强碰撞近似，不能把 $\varepsilon_0/\sigma$ 生硬解释为任意金属中电子惯性运动的真实完成时间。静电模型的使用应比较观测时间与材料响应时间，而不只是看电源是否标为直流。

